%MB9T5
\begin{vidu}%[3P5-3.2B][Loai2]
Trong thí nghiệm Young về giao thoa ánh sáng, nguồn sáng phát ra ánh sáng đơn sắc có bước sóng 450 nm. Khoảng cách giữa hai khe là 1 mm. Trên màn quan sát, khoảng cách giữa hai vân sáng liên tiếp là 0,72 mm.
\begin{itemize}
\item[1.]  Khoảng cách từ mặt phẳng chứa hai khe đến màn là $D=..........\ cm$
\item[2.] Vị trí vân sáng bậc 5 và vân tối thứ 4 là $x_{s5}=..........\ cm$ và  $x_{t4}=..........\ cm$
\item[3.] Xét điểm M trên màn có toạ độ 4,68 mm là vị trí vân sáng/tối bậc/thứ .......... 
\item[4.] Khoảng cách giữa 7 vân tối liên tiếp là $..........\ cm$
\item[5.] Khoảng cách giữa vân sáng bậc 3 và vân sáng bậc 5 ở hai phía so với vân sáng trung tâm là $..........\ cm$
\item[6.] Khoảng cách giữa vân sáng bậc 3 và vân tối thứ 5 ở cùng phía so với vân sáng trung tâm là $..........\ cm$
\end{itemize}
\loigiai{
\begin{itemize}
\item[1.] Công thức tính khoảng vân: $i=\dfrac{\lambda.D}{a}\Rightarrow D=\dfrac{ia}{\lambda}=\dfrac{0,72.10^{-3}.1.10^{-3}}{0,45.10^{-6}}=1,6\ m$.
\item[2.] Vị trí vân sáng bậc 5 ứng với \bth{k=\pm 5\Rightarrow x_{s(5)}=\pm 5i=\pm 3,6\ mm}.\\
Vị trí vân tối thứ 4 ứng với \bth{k=3, k=-4\Rightarrow x_{t(5)}=\pm 3,5i=\pm 2,52\ mm}.
\item[3.] Xét: $\dfrac{x_M}{i}=\dfrac{4,68}{0,72}=6,5$ nên M là vân tối thứ 7.
\item[4.]  Tìm khoảng cách giữa 7 vân tối liên tiếp ứng với 6 khoảng vân $6i=6.0,72=4,32\ mm$.
\item[5.] Do hai vân sáng khác phía: $\Delta x=3i+5i=8i=8.0,72=5,76\ mm$.
\item[6.] Do hai vân sáng cùng phía: $\Delta x=4,5i-3i=1,5i=1,5.0,72=1,08\ mm$.
\item
\item 
\end{itemize}
}
\end{vidu} \haidong{20}

\begin{vidu}%[3P5-3.2B][Loai2]
\nguon{QG - 2018} Trong thí nghiệm Young về giao thoa ánh sáng, nguồn sáng phát ra ánh sáng đơn sắc. Khoảng cách giữa hai khe là 1 mm, khoảng cách từ mặt phẳng chứa hai khe đến màn quan sát là 1,2 m. Trên màn, khoảng vân đo được là 0,6 mm. Bước sóng của ánh sáng trong thí nghiệm bằng
\choice
{600 nm}
{720 nm}
{480 nm}
{\True 500 nm}
\loigiai{
\begin{itemize}
\item Công thức tính bước sóng: $\lambda =\dfrac{ai}{D}=\dfrac{1.0,6}{1,2}=0,5\ \mu m=500\ nm$.
\end{itemize}
}
\end{vidu} \haidong{20}
\loai{Xác định vị trí vân}
\begin{vidu}%[3P5-3.2B][Loai3]
Trong thí nghiệm Young về giao thoa ánh sáng, các khe \bth{S_1,\ S_2} được chiếu bởi ánh sáng có bước sóng \bth{\lambda=0,54\microm}. Biết khoảng cách giữa hai khe là \bth{a=1,35\ mm}. Khoảng cách từ hai khe đến màn là \bth{D=1\ m}. Vị trí vân sáng bậc 5 và vân tối thứ 5 là
\choice
{\bth{x_{s(5)}=\pm 2\ mm} và \bth{x_{t(5)}=\pm 1,5\ mm}}
{\bth{x_{s(5)}=\pm 1,8\ mm} và \bth{x_{t(5)}=\pm 1,5\ mm}}
{\bth{x_{s(5)}=\pm 2\ mm} và \bth{x_{t(5)}=\pm 2,2\ mm}}
{\True \bth{x_{s(5)}=\pm 2\ mm} và \bth{x_{t(5)}=\pm 1,8\ mm}}
\loigiai{
\begin{itemize}
\item Đổi các đại lượng sang hệ đơn vị SI
\begin{align*}
\lambda=0,54\microm=0,54.10^{-6}\ m;\ D=1\ m ;\ a=1,35\ mm=1,35.10^{-3}\ m.
\end{align*}
\item Tính khoảng vân i: \bth{i=\dfrac{\lambda.D}{a}=\dfrac{0,54.10^{-6}.1}{1,35.10^{-3}}=0,4.10^{-3}\ m=0,4\ mm}.
\begin{itemize}
\item Vị trí vân sáng bậc 5 ứng với \bth{k=\pm 5\Rightarrow x_{s(5)}=\pm 5i=\pm 2\ mm}.
\item Vị trí vân tối thứ 5 ứng với \bth{k=4, k=-5\Rightarrow x_{t(5)}=\pm 4,5i=\pm 1,8\ mm}.
\end{itemize}
\end{itemize}
}
\end{vidu} \haidong{20}

\loai{Từ vị trí vân tìm bước sóng}
\begin{vidu}%[3P5-3.2B][Loai4]
Cho hai nguồn sáng kết hợp $S_1$ và $S_2$ cách nhau một khoảng $a=5\ mm$ và cách đều một màn E một khoảng $D=2\ m$. Quan sát vân giao thoa trên màn, người ta thấy khoảng cách từ vân sáng bậc 5 đến vân trung tâm là 1,5\ mm. Bước sóng $\lambda$ của nguồn sáng là
\choice
{$0,5\ \mu m$}
{$0,55\ \mu m$}
{$0,60 \ \mu m$}
{\True $0,75\ \mu m$}
\loigiai{
\begin{itemize}
\item Khoảng cách từ vân sáng bậc 5 đến vân trung tâm bằng 5 khoảng vân nên
\begin{align*}
i=\dfrac{1,5}{5}=0,3\ mm\Rightarrow \lambda =\dfrac{ia}{D}=\dfrac{0,3. 5}{2}=0,75\ \mu m.
\end{align*}
\item 
\end{itemize}
}
\end{vidu} \haidong{20}
\loai{Xác định M thuộc vân sáng hay vân tối}
\begin{vidu}%[3P5-3.2B][Loai5]
Trong thí nghiệm giao thoa Young, khoảng cách hai khe là 1,2\ mm, khoảng cách giữa mặt phẳng chứa hai khe và màn ảnh là 2 m. Người ta chiếu vào khe Young bằng ánh sáng đơn sắc có bước sóng $ 0,6\ \mu m $. Xét tại hai điểm M và N trên màn có tọa độ lần lượt là 6\ mm và 15,5\ mm là vị trí vân sáng hay vân tối?
\choice
{M sáng bậc 2; N tối thứ 16}
{\True M sáng bậc 6; N tối thứ 16}
{M sáng bậc 2; N tối thứ 9}
{M tối bậc 2; N tối thứ 9}
\loigiai{
\begin{itemize}
\item Ta có: $i=\dfrac{\lambda D}{a}=\dfrac{{0,6.10}^{-6}.2}{{1,2.10}^{-3}}=1\ mm\Rightarrow \left\{\begin{aligned}& \dfrac{x_m}{i}=6\Rightarrow M\ \text{là vân sáng bậc 6}\\
& \dfrac{x}{i}=15,5\Rightarrow N\ \text{là vân tối thứ 16}
\end{aligned}\right.$
\end{itemize}
}
\end{vidu} \haidong{20}
\loai{Cho khoảng cách giữa n vân liên tiếp. Tính các đại lượng}
\begin{vidu} %[3P5-3.2K][Loai7]
\nguon{QG - 2020} Trong thí nghiệm Young về giao thoa ánh sáng đơn sắc, khoảng cách giữa 4 vân sáng liên tiếp trên màn quan sát là 2,4 mm. Khoảng vân trên màn là
\choice
{1,6 mm}
{1,2 mm}
{0,6 mm}
{\True 0,8 mm}
\loigiai{
\begin{itemize}
\item 4 vân sáng: $3i = 2,4 \Rightarrow i = 0,8\ mm$.
\end{itemize}
}
\end{vidu} \haidong{20}
\loai{Xác định khoảng cách giữa hai vân sáng}
\begin{vidu}%[3P5-3.2K][Loai8]
\nguon{MH - 2019} Tiến hành thí nghiệm Young về giao thoa ánh sáng với ánh sáng đơn sắc có bước sóng $0,6\ \mu m$. Khoảng cách giữa hai khe là 0,3 mm, khoảng cách từ mặt phẳng chứa hai khe đến màn quan sát là 2 m. Trên màn, khoảng cách giữa vân sáng bậc 3 và vân sáng bậc 5 ở hai phía so với vân sáng trung tâm là
\choice
{8 mm}
{\True 32 mm}
{20 mm}
{12 mm}
\loigiai{
\begin{itemize}
\item Do hai vân sáng khác phía: $\Delta x=\dfrac{\lambda D}{a}\left(k_5+k_3\right)=\dfrac{0,6.10^{-6}.2}{0,3.10^{-3}}(5+3)=0,032\ m=32\ mm$.
\end{itemize}
}
\end{vidu} \haidong{20}
\loai{Khoảng cách giữa hai vân tối hoặc giữa vân tối với vân sáng}
\begin{vidu}%[3P5-3.2K][Loai9]
Trong thí nghiệm Young về giao thoa ánh sáng, hai khe được chiếu bằng ánh sáng đơn sắc có bước sóng 0,6\microm. Khoảng cách giữa hai khe sáng là 1\ mm, khoảng cách từ mặt phẳng chứa hai khe đến màn quan sát là 1,5\ m. Trên màn quan sát, khoảng cách từ vân sáng bậc 2 đến vân tối thứ 10 khác phía so với vân trung tâm cách nhau một đoạn là
\choice
{6,45 mm}
{\True 10,35 mm}
{10,9 mm}
{8,8 mm}
\loigiai{
\begin{itemize}
\item Khoảng vân giao thoa trên màn là \bth{i=\dfrac{\lambda.D}{a}=\dfrac{0,6.1,5}{1}=0,9\ mm}.
\begin{itemize}
\item Vị trí vân sáng bậc 2: \bth{x_{s(2)}=2i=2.0,9=1,8\ mm}
\item Vị trí vân tối thứ 10: \bth{x_{t(10)}=\pbrace{9+\dfrac{1}{2}}i=\pbrace{9+\dfrac{1}{2}}.0,9=8,55\ mm}
\end{itemize}
\item Vì hai vân khác phía nên khoảng cách giữa hai vân là
\begin{align*}
\Delta x=x_{s(2)}+x_{t(10)}=1,8+8,55=10,35\ mm.
\end{align*}
\end{itemize}
}
\end{vidu} \haidong{20}
\loai{Từ khoảng cách giữa hai vân tìm các đại lượng}
\begin{vidu}%[3P5-3.2K][Loai10]
Trong thí nghiệm Young về giao thoa ánh sáng, khoảng cách giữa hai khe là 1,5 mm, khoảng cách từ hai khe đến màn là 3 m, người ta đo được khoảng cách giữa vân sáng bậc 2 đến vân sáng bậc 5 ở cùng phía với nhau so với vân sáng trung tâm là 3 mm. Bước sóng của ánh sáng dùng trong thí nghiệm là
\choice
{0,2 $\mu m $}
{0,4 $\mu m $}
{\True 0,5 $\mu m $}
{0,6 $\mu m $}
\loigiai{
\begin{itemize}
\item Khoảng cách giữa vân sáng bậc 2 đến vân sáng bậc 5 cùng phía với nhau so với vân sáng trung tâm bằng:
$5i-2i=3i=3\ mm\Rightarrow i=1mm\Rightarrow \lambda =\dfrac{ia}{D}=\dfrac{1. 1,5}{3}=0,5\ \mu m $.
\end{itemize}
}
\end{vidu} \haidong{20}
\loai{Hiệu đường đi}
\begin{vidu}%[3P5-3.2K][Loai11]
Trong thí nghiệm Young về giao thoa ánh sáng, hai khe được chiếu bằng ánh sáng đơn sắc có bước sóng $\lambda $. Nếu tại điểm M trên màn quan sát có vân tối thứ 3 (tính từ vân sáng trung tâm) thì hiệu đường đi của ánh sáng từ hai khe $S_1,\ S_2$ đến M có độ lớn bằng
\choice
{\True $2,5\lambda $}
{$2\lambda $}
{$3\lambda $}
{$1,5\lambda $}
\loigiai{\begin{itemize}
\item Vân tối thứ 3 ứng với $\left[\begin{aligned}
&{k=2} \\
&{k=-3} \\
\end{aligned}\right.$
\item Hiệu đường đi của tia sáng tới hai khe: $d_2-d_1=\left({2k+1}\right)\dfrac{\lambda}{2}=\left({2.2+1}\right)\dfrac{\lambda}{2}=2,5\lambda $.
\end{itemize}}
\end{vidu} \haidong{20}
\loai{Điều kiện kẹp bước sóng}
\begin{vidu}%[3P5-3.2K][Loai12]
Trong thí nghiệm Young về giao thoa ánh sáng, khoảng cách giữa hai khe là 2 mm, khoảng cách từ hai khe đến màn là 2 m, ánh sáng đơn sắc dùng trong thí nghiệm có bước sóng trong khoảng từ $0,40\ \mu m$ đến $0,76\ \mu m$. Tại vị trí cách vân sáng trung tâm 1,56 mm là một vân sáng. Bước sóng của ánh sáng dùng trong thí nghiệm là
\choice
{$\lambda =0,42\ \mu m$}
{$\lambda =0,62\ \mu m$}
{\True $\lambda =0,52\ \mu m$}
{$\lambda =0,72\ \mu m$}
\loigiai{\begin{itemize}
\item Bước sóng của bức xạ cho vân sáng tại vị trí x: $x=k.\dfrac{\lambda D}{a}\Rightarrow \lambda =\dfrac{ax}{k.D}=\dfrac{2.1,56}{k.2}=\dfrac{1,56}{k}\ \mu m$.
\item Cho $\lambda $ vào điều kiện bước sóng ta có: $0,4\le \dfrac{1,56}{k}\le 0,76\Rightarrow 2,05\le k\le 3,9\Rightarrow k=3$.
\item Bước sóng của ánh sáng dùng trong thí nghiệm là: $\lambda =\dfrac{1,56}{3}=0,52\ \mu m$.
\end{itemize}}
\end{vidu} \haidong{20}
\loai{Cho khoảng bước sóng và hai đoạn. Tìm bước sóng}
 
\begin{vidu} %[3P5-3.2K][Loai14]
\nguon{QG - 2019} Tiến hành thí nghiệm Young về giao thoa ánh sáng, nguồn sáng phát ra ánh sáng đơn sắc có bước sóng $\lambda $ ($380\ nm < \lambda  < 760\ nm$). Khoảng cách giữa hai khe là 1 mm, khoảng cách từ mặt phẳng chứa hai khe đến màn quan sát là 1 m. Trên màn hai điểm A và B là vị trí vân sáng đối xứng với nhau qua vân trung tâm, C cũng là vị trí vân sáng. Biết A, B, C cùng nằm trên một đường thẳng vuông góc với các vân giao thoa, $AB = 6,6$ mm; $BC = 4,4$ mm. Giá trị của $\lambda $ bằng
\choice
{\True 550 nm}
{450 nm}
{750 nm}
{650 nm}
\loigiai{
\begin{itemize}
\item Ta có $OA=OB=3,3$\ mm. 
\item Tại A, B, C đều là các vân sáng nên ta có:\\ $OC=CB-OB=1,1\ mm=11.10^5\ nm=k\dfrac{\lambda D}{a}=1000.k\lambda $ với k nguyên thì chỉ có $\lambda =550\ nm$ thỏa mãn.
\end{itemize}}
\end{vidu} \haidong{20}